\documentclass[book.tex]{subfiles}
\begin{document}

\chapter{ANI Potentials}
\label{chap:ani}
\noindent\rule{11cm}{0.4pt} \\
\textbf{Prerequisites:}  \autoref{chap:molecular_hamiltonian},   \\
\textbf{Difficulty Level:} ** \\
\noindent\rule{11cm}{0.4pt}
\vspace{5mm}

ANI-1 \cite{smith2017anidata}, \cite{smith2019approaching}, \cite{smith2020ani}, \cite{gao2020torchani}, \cite{devereux2020extending} is a deep neural network capable of learning accurate and transferable potentials for organic molecules. It is based on the Behler-Parinello symmetry function method \cite{behler2007generalized} with additional changes enabling it to learn different potentials for different atom types. A feature vector, a series of symmetry functions, is built for each atom in the molecule based on its atom type and interaction with other atoms. Then the feature vectors are fed into different neural network potentials (depending on atom types) to generate predictions of properties. %This model is first introduced by Smith et al.\cite{smith2017aninn} In their original article, the model is trained on 58k small molecules with 8 or less heavy atoms, each with multiple poses and potentials. Training set in total has 17.2 million data points, which is far bigger than qm8 or qm9 in our collection. Since we only have molecules in their most stable configuration, we cannot expect similar level of accuracy. Further comparison and benchmarking with similar size of training set is left to future work.

\section{Mathematical Formulation}

The base structure of the ANI potential closely follows the Behler-Parinello framework. Atomic energies are estimated as the sum of local atomic contributions

\begin{align}
E &= \sum_i E_i
\end{align}

The local atomic environment is again computed with function $f_C$
\begin{align}
f_C(R_{ij}) &= \begin{cases}
0.5 \left [ \cos \left ( \frac{\pi R_{ij}}{R_c} +  1 \right ) \right ] & R_{ij} < R_C \\
0.0 & R_{ij} \geq R_C
\end{cases}
\end{align}

The atomic symmetry functions are computed as before
\begin{align}
G^R_m &= \sum_{j \neq i} e^{-\eta (R_{ij} - R_s)^2 } f_C(R_{ij})
\end{align}
The index $m$ here ranges over different choices of $\eta$ and $R_s$. 


Angular contribution functions are modified by the addition of an offset parameter $\theta_s$ and the addition of $R_s$ to the exponential term.
\begin{align}
G^{A_{mod}}_m &= 2^{1-\zeta} \sum_{j,k \neq i} (1 + \cos(\theta_{ijk} - \theta_s))^\zeta \exp \left [ \eta \left (\frac{R_{ij} + R_{ik}}{2} - R_s \right )^2 \right ] f_C(R_{ij}) f_C(R_{ik}) 
\end{align}
The index $m$ here ranges over $\zeta, \theta_s$, $\eta$, and $R_s$. 

These modified symmetry functions are used to compute an atomic environment vector (AEV) at each atom

\begin{align}
G^X_i &= \{G_1, \dotsc, G_M\}
\end{align}

The number of different features $M$ is a hyperparameter of the system. The original paper chooses a total of $768$ different features by setting different values of the hyperparameters.

\section{Normal Mode Sampling}

The ANI functions can be used to sample molecular conformations. Normal mode coordinates have to be computed for the potential energy surface and a displace is computed along each coordinate.

\section{Training}

The ANI function is trained using an exponential loss function that compares against energies computed with density functional theory.

\begin{align}
\mathcal{L}(E^{ANI}) &= \tau \exp \left ( \frac{1}{\tau} \sum_j (E^{ANI}_j E^{DFT}_j)^2 \right )
\end{align}
\end{document}
